\chapter{Basis, Dimension, and Linear Transformations}

An anomaly detector does not read raw network traffic. It summarizes it first. Four
counters, say failed connections, outbound packets, session length, and the number of
unique ports, are compressed into three features by the linear map
$F\colon\mathbb{R}^{4}\to\mathbb{R}^{3}$.
The decision threshold only looks at $F(\mathbf{x})$.

The decisive question: is there $\boldsymbol{\delta}\neq\mathbf{0}$ with
$F(\mathbf{x}+\boldsymbol{\delta})=F(\mathbf{x})$? Because $F$ is linear, that condition is
equivalent to
$F(\boldsymbol{\delta})=\mathbf{0}$. The answer is a subspace of dimension
$4-\operatorname{rank}F$; if $F$ has rank three, there is a whole line of blind directions,
and an attacker can shift all four counters along that line without shifting a single one of
the features.

To answer it we need three things Chapter 6 has not yet given: the list of coordinates with
respect to a basis, the rule for moving that list when the basis changes, and the translation
of a linear map into a matrix so that the kernel can be computed by elimination. Chapter 6
already established that a basis exists and that its size is unique; its axioms are not
repeated here.

\begin{untukanak}
Imagine a bookcase with many shelves. Every book has a fixed place. If we rearrange the
bookcase, the books do not change, but each book's address does. To say where a book is, we
must name the arrangement currently in use. A friend who uses another arrangement will name a
different address for the same book,
and the two of you are both right.

There is also a special shelf behind the bookcase. A book moved there cannot be seen from the
front. As long as books only move within that shelf, the list of books in front does not
change at all.

Have you ever moved something to a place nobody inspects? We are going to find out why the
list of addresses can change while the contents of
the bookcase stay the same, why that hidden shelf always has a definite size, and why some
changes do not shift the summary at all.
\end{untukanak}

\section{Coordinates and transition matrices}
\tingkat{intermediate}

Vectors in $\mathbb{R}^{3}$ are already lists of three numbers; vectors in $P_{2}$ or
$M_{2\times2}(\mathbb{R})$ are not. So that both can be computed in the same way, we
need a translation that turns every vector into a list of numbers. That translation depends
on the basis, and that is where the difficulty of this section begins.

Uniqueness we already have: the uniqueness lemma for coordinates in Chapter 6 states that the
expression of a vector in a basis is always unique. Therefore the following definition is
legitimate without choosing among several expressions.

\kamus{coordinate vector}{koordinat}{the address of a book in the bookcase: its contents stay
fixed, its address depends on how the shelves are arranged}

\begin{definisi}[Coordinate vector]
Let $B=(\mathbf{b}_{1},\dots,\mathbf{b}_{n})$ be an ordered basis of the vector space
$V$
over $\mathbb{R}$. For every $\mathbf{v}\in V$ there are unique scalars $a_{1},\dots,a_{n}$
with $\mathbf{v}=a_{1}\mathbf{b}_{1}+\dots+a_{n}\mathbf{b}_{n}$. The column
\[
  [\mathbf{v}]_{B}=(a_{1},\dots,a_{n})^{\top}
\]
is called the \emph{coordinate vector} of $\mathbf{v}$ with respect to $B$.\index{coordinate}
\end{definisi}

The word ``ordered'' is not decoration. If two basis vectors are exchanged, two coordinate
entries are exchanged with them while the vector stays put. A basis set alone does not
determine the coordinate list;
its order does as well.

One warning: $[\mathbf{v}]_{B}$ is not $\mathbf{v}$. Those numbers live in
$\mathbb{R}^{n}$, whereas $\mathbf{v}$ lives in $V$. For $M_{2\times2}(\mathbb{R})$
the difference is palpable: the vector holds four numbers, and its coordinates are four
numbers too, but arranged as a column that depends on the basis.

The first important property of this translation: it respects both operations of the vector
space.

\begin{proposisi}
For every basis $B$ of a space $V$ of dimension $n$, the map
$\varphi_{B}\colon V\to\mathbb{R}^{n}$ with $\varphi_{B}(\mathbf{v})=[\mathbf{v}]_{B}$
is linear, injective, and surjective.
\end{proposisi}

\begin{proof}
Take $\mathbf{v}=\sum_{i}a_{i}\mathbf{b}_{i}$ and $\mathbf{w}=\sum_{i}c_{i}\mathbf{b}_{i}$.
Then $\mathbf{v}+\mathbf{w}=\sum_{i}(a_{i}+c_{i})\mathbf{b}_{i}$. Uniqueness of coordinates
forces $[\mathbf{v}+\mathbf{w}]_{B}=(a_{i}+c_{i})=[\mathbf{v}]_{B}+[\mathbf{w}]_{B}$.
In the same way, $[c\mathbf{v}]_{B}=c[\mathbf{v}]_{B}$.

For injectivity, suppose $[\mathbf{v}]_{B}=\mathbf{0}$. Then all its coefficients are zero,
so $\mathbf{v}=\sum_{i}0\cdot\mathbf{b}_{i}=\mathbf{0}$. For surjectivity, take
any $(a_{1},\dots,a_{n})\in\mathbb{R}^{n}$ and form
$\mathbf{v}=\sum_{i}a_{i}\mathbf{b}_{i}$; its coordinates are exactly that list.
\end{proof}

The map $\varphi_{B}$ turns a question about $V$ into a question about columns in
$\mathbb{R}^{n}$, then turns it back; that transfer depends on $B$. Coordinates
are written as columns so that matrices act from the left: coordinates in, multiplied by the
matrix, coordinates out. With row coordinates, every formula would have to be transposed.

\begin{bacakeras}
Two voices stand in front of the class.

``How long is this table?''
``Three.''
``Three what?''
``Three rulers. Or three spans. Or three paces. Choose the measuring instrument first.''

It is the same with a vector space. A vector has no address until we choose a basis. The basis
is
the ruler: it supplies a scale, a direction, and an order. Change the ruler and you change
the numbers. The object being measured does not change with them.

One ruler, one number. Two rulers, two numbers. Without a ruler, no numbers
at all.

Be careful with a bad ruler: two of its scale marks are almost parallel. Numbers still come
out,
but the numbers wobble, and two people measuring patiently will obtain different
numbers.

Repeat together: a basis is not the object, a basis is the instrument. What changes is the
address, what stays is the content.
\end{bacakeras}

Now two bases are used on the same space. Every vector has two coordinate lists,
and we want a formula that turns the first list into the second. The matrix that
performs that task is built from the coordinates of the old basis vectors in the new basis.

\begin{definisi}[Transition matrix]
Let $B=(\mathbf{b}_{1},\dots,\mathbf{b}_{n})$ and
$C=(\mathbf{c}_{1},\dots,\mathbf{c}_{n})$
be two bases of $V$. The \emph{transition matrix} from $B$ to $C$ is the $n\times n$ matrix
\[
P_{C\leftarrow
B}=\Bigl(\,[\mathbf{b}_{1}]_{C}\;\Big|\;\cdots\;\Big|\;[\mathbf{b}_{n}]_{C}\,\Bigr),
\]
that is, the matrix whose $j$-th column is $[\mathbf{b}_{j}]_{C}$.\index{transition matrix}
\end{definisi}

What still needs proof: multiplying the old coordinate list by this matrix produces the new
coordinate
list.

\begin{teorema}
For every $\mathbf{v}\in V$ we have $[\mathbf{v}]_{C}=P_{C\leftarrow B}[\mathbf{v}]_{B}$.
\end{teorema}

\begin{proof}
Write $\mathbf{v}=\sum_{j}a_{j}\mathbf{b}_{j}$, so that
$[\mathbf{v}]_{B}=(a_{1},\dots,a_{n})^{\top}$.
Linearity of $\varphi_{C}$ gives
$[\mathbf{v}]_{C}=\sum_{j}a_{j}[\mathbf{b}_{j}]_{C}$. That right-hand side is exactly the
matrix
product $P_{C\leftarrow B}$ with the column $(a_{j})$, by the definition of matrix-vector
multiplication.
\end{proof}

A corollary that makes the matrix usable in both directions.

\begin{akibat}
The transition matrix is always invertible, and $P_{B\leftarrow C}=P_{C\leftarrow B}^{-1}$.
Moreover
$P_{C\leftarrow B}=P_{C\leftarrow D}\,P_{D\leftarrow B}$ for three bases $B,C,D$ of $V$.
\end{akibat}

\begin{proof}
The previous theorem gives $P_{B\leftarrow C}P_{C\leftarrow B}[\mathbf{v}]_{B}=[\mathbf{v}]_{B}$
for all $\mathbf{v}$. Because $\varphi_{B}$ is surjective, that product equals the identity
matrix. The same argument applies in the opposite order, so the two matrices are mutually
inverse. For the chain rule, apply the theorem twice:
$[\mathbf{v}]_{C}=P_{C\leftarrow D}[\mathbf{v}]_{D}=P_{C\leftarrow D}P_{D\leftarrow B}[\mathbf{v}]_{B}$,
then compare with $[\mathbf{v}]_{C}=P_{C\leftarrow B}[\mathbf{v}]_{B}$ at $\mathbf{v}=\mathbf{b}_{j}$.
\end{proof}

Note that $P_{C\leftarrow B}$ and $P_{B\leftarrow C}$ are not transposes of each other; they are mutual inverses, and that is a different thing. The direction of the subscript is read from right to left, like the composition of maps: take the $j$-th basis vector of $B$, then ask for its address according to $C$. When $B$ is the standard basis, the columns of $P_{B\leftarrow C}$ are the vectors of $C$ themselves.

Take $B=(\mathbf{e}_{1},\mathbf{e}_{2},\mathbf{e}_{3})$
and $C=(\mathbf{b}'_{1},\mathbf{b}'_{2},\mathbf{b}'_{3})$ with
\[
  \mathbf{b}'_{1}=(1,1,0),\qquad \mathbf{b}'_{2}=(1,0,1),\qquad \mathbf{b}'_{3}=(0,1,1).
\]
Because $B$ is standard, the matrix $P_{B\leftarrow C}$ contains those three vectors as columns:
\[
  P_{B\leftarrow C}=\begin{pmatrix}1&1&0\\ 1&0&1\\ 0&1&1\end{pmatrix},
  \qquad \det P_{B\leftarrow C}=-2\neq0 .
\]
Its inverse is
\[
  P_{C\leftarrow B}=\frac{1}{2}\begin{pmatrix}1&1&-1\\ 1&-1&1\\ -1&1&1\end{pmatrix}.
\]
For $\mathbf{v}=(3,8,11)$ we obtain
\[
  [\mathbf{v}]_{C}=P_{C\leftarrow B}[\mathbf{v}]_{B}
  =\frac{1}{2}\begin{pmatrix}3+8-11\\ 3-8+11\\ -3+8+11\end{pmatrix}
  =\begin{pmatrix}0\\ 3\\ 8\end{pmatrix}.
\]
The check is immediate: $3\mathbf{b}'_{2}+8\mathbf{b}'_{3}=(3,8,11)$. Conversely
$\mathbf{w}=(1,0,0)$ has coordinates $(1/2,\,1/2,\,-1/2)$; a fractional entry is natural,
because
coordinates are nothing but an address.

\section{Linear transformations, kernel, and range}
\tingkat{basic}

Chapter 6 discussed spaces and subspaces, but not maps between them. Yet the main work of
linear algebra is moving vectors: rotating, projecting, differentiating, summarizing.

A rule that respects addition and scalar multiplication is called a linear transformation:
adding first and then mapping gives the same result as mapping first and then adding.

\begin{definisi}[Linear transformation]
A map $T\colon V\to W$ between two vector spaces over $\mathbb{R}$ is called a
\emph{linear transformation} if $T(\mathbf{u}+\mathbf{v})=T(\mathbf{u})+T(\mathbf{v})$ and
$T(c\mathbf{v})=cT(\mathbf{v})$ for all $\mathbf{u},\mathbf{v}\in V$ and $c\in\mathbb{R}$.
\end{definisi}

The two conditions are equivalent to a single condition,
$T(c\mathbf{u}+\mathbf{v})=cT(\mathbf{u})+T(\mathbf{v})$. The first consequence,
$T(\mathbf{0})=\mathbf{0}$, is used to reject linearity quickly: the map
$T(\mathbf{x})=\mathbf{x}+\mathbf{a}$ with $\mathbf{a}\neq\mathbf{0}$ fails at once.

Conversely $T(\mathbf{0})=\mathbf{0}$ is not enough. The map $T(x,y)=(x^{2},y)$ satisfies it,
but
$T(2,0)=(4,0)$ whereas $2T(1,0)=(2,0)$. Squaring breaks linearity.\index{linear
transformation}

A $90^{\circ}$ rotation, $R(x,y)=(-y,x)$, is linear: $R(3,1)=(-1,3)$. The derivative $D\colon P_{3}\to P_{2}$, $D(p)=p'$, is linear because $(p+q)'=p'+q'$. The map $T(x,y)=(x+y,\,2x+2y)$ is linear too, but all of $\mathbb{R}^{2}$ falls onto a single line; it is linear and simultaneously destroys information.

Take
$A=\begin{pmatrix}2&1\\1&2\end{pmatrix}$ and $T(\mathbf{x})=A\mathbf{x}$. The vector $(1,0)$ goes
to $(2,1)$ and $(0,1)$ to $(1,2)$, so the unit square becomes a parallelogram with
vertices
$(0,0)$, $(2,1)$, $(3,3)$, $(1,2)$ of area $3$, since $|2\cdot2-1\cdot1|=3$.
Figure~\ref{fig:kisi-07} shows the grid before and after the map: lines stay lines, the grid
stays a grid, but distances and areas are no longer the same.

\begin{figure}[htbp]
\centering
\begin{tikzpicture}[scale=0.38,>=stealth,font=\scriptsize]
  \foreach \k in {0,1,2,3}{
    \draw[layergrey!45] (\k,0)--(\k,3);
    \draw[layergrey!45] (0,\k)--(3,\k);
    \draw[layerteal,thick] (2*\k,\k)--(2*\k+3,\k+6);
    \draw[layerteal,thick] (\k,2*\k)--(\k+6,2*\k+3);}
  \fill[accent!15] (0,0)--(1,0)--(1,1)--(0,1)--cycle;
  \draw[accent,thick] (0,0)--(1,0)--(1,1)--(0,1)--cycle;
  \fill[layerorange!25] (0,0)--(2,1)--(3,3)--(1,2)--cycle;
  \draw[layerorange,very thick] (0,0)--(2,1)--(3,3)--(1,2)--cycle;
  \draw[->,ink,thick] (-0.6,0)--(9.8,0) node[below right]{$x$};
  \draw[->,ink,thick] (0,-0.6)--(0,9.8) node[left]{$y$};
  \foreach \x/\y in {2/1,1/2,3/3}{
    \fill[ink] (\x,\y) circle (1.8pt);
    \node[ink,anchor=north west] at (\x+0.05,\y-0.05) {$(\x,\y)$};}
  \node[layergrey,anchor=north west] at (0,-1.0) {thin: original grid};
  \node[layerteal,anchor=north west] at (4.0,-1.0) {thick: grid after mapping};
\end{tikzpicture}
\caption{The unit grid and its image under $T(\mathbf{x})=A\mathbf{x}$ with
$A=\begin{pmatrix}2&1\\1&2\end{pmatrix}$. One cell becomes a parallelogram with
vertices $(0,0)$, $(2,1)$, $(3,3)$, and $(1,2)$, so its area grows from $1$ to $3$.
The old grid is drawn thin, the new grid thick; both remain grids.}
\label{fig:kisi-07}
\end{figure}

\begin{tangan}
Take a sheet of graph paper, a pencil, and a ruler. Draw a horizontal axis and a vertical
axis,
then mark their intersection as the center. Draw a small house: a square five cells wide, a
triangular
roof, a door one cell wide.

Now translate. Move every point of the house three cells to the right and two cells up.
Do not guess; count. A point that was one cell right and one cell up
is now
four cells right and three cells up. Redraw the house. Its shape and
size are the same, only its place has moved.

Now rotate. Turn the paper a quarter turn clockwise, then redraw the same
house. The length of each side stays, the direction of each side changes.

What should be visible: under translation, all directions are preserved but the center moves with the figure; under rotation, the center stays and all directions change. Because a translation moves the center, a translation is
not
a linear transformation. A rotation is.
\end{tangan}

The good news: to know a linear map, it is enough to know its values on the basis vectors.

\begin{proposisi}
Let $B=(\mathbf{b}_{1},\dots,\mathbf{b}_{n})$ be a basis of $V$ and let
$\mathbf{w}_{1},\dots,\mathbf{w}_{n}\in W$ be arbitrary. There is exactly one linear map
$T\colon V\to W$ with $T(\mathbf{b}_{j})=\mathbf{w}_{j}$ for every $j$.
\end{proposisi}

\begin{proof}
For existence, define
$T\bigl(\sum_{j}a_{j}\mathbf{b}_{j}\bigr)=\sum_{j}a_{j}\mathbf{w}_{j}$.
This definition does not depend on the expression, because coordinates with respect to a basis
are unique. For linearity, use the linearity of coordinates from the proposition in the
previous section:
$T(\mathbf{u}+\mathbf{v})=\sum_{j}(a_{j}+c_{j})\mathbf{w}_{j}=T(\mathbf{u})+T(\mathbf{v})$,
and similarly for scalar multiplication.

For uniqueness, suppose $S$ is linear and $S(\mathbf{b}_{j})=\mathbf{w}_{j}$. For
$\mathbf{v}=\sum_{j}a_{j}\mathbf{b}_{j}$ we have
$S(\mathbf{v})=\sum_{j}a_{j}S(\mathbf{b}_{j})=\sum_{j}a_{j}\mathbf{w}_{j}=T(\mathbf{v})$.
\end{proof}

The practical corollary: $n$ vectors in $W$ determine one linear map from an $n$-dimensional
$V$.
When $W=\mathbb{R}^{m}$, that data is nothing other than an $m\times n$ matrix.

Every linear map carries two sets worth inspecting: the vectors mapped to zero,
and the values that are actually attained.

\kamus{kernel or null space}{inti}{the list of changes that do not shift the
summary at all}

\begin{definisi}[Kernel, range, rank, and nullity]
For a linear $T\colon V\to W$, the \emph{kernel} is
$\ker T=\{\mathbf{v}\in V: T(\mathbf{v})=\mathbf{0}\}$ and the \emph{range} is
$\operatorname{im}T=\{T(\mathbf{v}):\mathbf{v}\in V\}$. When $V$ is finite-dimensional,
$\operatorname{nullity}(T)=\dim\ker T$ and $\operatorname{rank}(T)=\dim\operatorname{im}T$.
\end{definisi}

Both are subspaces, not merely sets.\index{kernel}\index{range}

\begin{lemma}
$\ker T$ is a subspace of $V$ and $\operatorname{im}T$ is a subspace of $W$.
\end{lemma}

\begin{proof}
Both contain $\mathbf{0}$ because $T(\mathbf{0})=\mathbf{0}$. If
$\mathbf{u},\mathbf{v}\in\ker T$
and $c\in\mathbb{R}$, then $T(\mathbf{u}+\mathbf{v})=\mathbf{0}$ and
$T(c\mathbf{u})=\mathbf{0}$,
so $\ker T$ is closed. If $\mathbf{u}=T(\mathbf{x})$ and $\mathbf{v}=T(\mathbf{y})$, then
$\mathbf{u}+\mathbf{v}=T(\mathbf{x}+\mathbf{y})$ and $c\mathbf{u}=T(c\mathbf{x})$ belong to
$\operatorname{im}T$.
\end{proof}

The kernel measures the failure of injectivity, and that is all that needs checking.

\begin{proposisi}
$T$ is injective if and only if $\ker T=\{\mathbf{0}\}$.
\end{proposisi}

\begin{proof}
If $T$ is injective, only $\mathbf{0}$ maps to $\mathbf{0}$. Conversely, suppose
$\ker T=\{\mathbf{0}\}$ and $T(\mathbf{u})=T(\mathbf{v})$. Then $T(\mathbf{u}-\mathbf{v})=\mathbf{0}$,
so $\mathbf{u}-\mathbf{v}\in\ker T$, hence $\mathbf{u}=\mathbf{v}$.
\end{proof}

Rank counts the directions that survive, nullity counts the directions that vanish. Together
they divide the dimension of the source space.

\kamus{rank}{rank}{the number of independent directions that still survive after the map}

A small example shows both fates. Take
$B=\begin{pmatrix}1&2\\2&4\end{pmatrix}$. Then $B\mathbf{e}_{1}=(1,2)$ and
$B\mathbf{e}_{2}=(2,4)$, two collinear vectors, so $\operatorname{rank}(B)=1$. The direction $(1,2)$
survives and stretches: $B(1,2)=(5,10)$. The direction $(2,-1)$ vanishes: $B(2,-1)=(0,0)$.
Figure~\ref{fig:arah} draws both, and rank-nullity closes the count: $2=1+1$.

\begin{figure}[htbp]
\centering
\begin{tikzpicture}[scale=0.60,>=stealth,font=\scriptsize]
  \draw[layerjade!60,thin] (-0.7,-1.4)--(2.8,5.6);
  \node[layerjade,anchor=west,rotate=63] at (2.0,4.3) {line in direction $(1,2)$};
  \draw[->,ink,thick] (-1.0,0)--(3.3,0) node[below right]{$x$};
  \draw[->,ink,thick] (0,-2.0)--(0,5.0) node[left]{$y$};
  \draw[->,layergrey,thick] (0,0)--(1,0) node[below]{$\mathbf{e}_{1}$};
  \draw[->,layergrey,thick] (0,0)--(0,1) node[left]{$\mathbf{e}_{2}$};
  \draw[->,layerorange,very thick] (0,0)--(1,2) node[above left]{$B\mathbf{e}_{1}=(1,2)$};
  \draw[->,layerbrown,very thick] (0,0)--(2,4) node[right]{$B\mathbf{e}_{2}=(2,4)$};
  \draw[->,layerplum,thick,dashed] (0,0)--(1.8,-0.9) node[below right]{$(2,-1)$};
  \node[layerplum,anchor=north] at (1.5,-1.6) {collapses to $\mathbf{0}$};
  \fill[ink] (0,0) circle (1.6pt) node[below left,ink]{$\mathbf{0}$};
\end{tikzpicture}
\caption{Two directions under $B=\begin{pmatrix}1&2\\2&4\end{pmatrix}$. The direction $(1,2)$
survives and stretches into $(5,10)$; the direction $(2,-1)$ collapses to $\mathbf{0}$. The
images of the two basis vectors,
$B\mathbf{e}_{1}=(1,2)$ and $B\mathbf{e}_{2}=(2,4)$, lie on a single line, so
$\operatorname{rank}(B)=1$ and $\operatorname{nullity}(B)=1$.}
\label{fig:arah}
\end{figure}

Now the central theorem of this chapter. Its content: no dimension is lost along the way.

\begin{teorema}[Rank-nullity]
If $V$ is finite-dimensional and $T\colon V\to W$ is linear, then
\[
  \dim V=\dim\ker T+\dim\operatorname{im}T .
\]
\end{teorema}

\begin{proof}
Let $\dim V=n$ and $\dim\ker T=k$. Take a basis
$(\mathbf{u}_{1},\dots,\mathbf{u}_{k})$ of $\ker T$. By the basis extension corollary in Chapter 6,
that basis can be completed to a basis of $V$ by adding $n-k$ vectors
$\mathbf{w}_{1},\dots,\mathbf{w}_{n-k}$.

First, the set $\{T(\mathbf{w}_{1}),\dots,T(\mathbf{w}_{n-k})\}$ spans
$\operatorname{im}T$. For an arbitrary $T(\mathbf{v})$, write
$\mathbf{v}=\sum_{i}\alpha_{i}\mathbf{u}_{i}+\sum_{j}\beta_{j}\mathbf{w}_{j}$. The terms
$\mathbf{u}_{i}$ vanish because $T(\mathbf{u}_{i})=\mathbf{0}$, so
$T(\mathbf{v})=\sum_{j}\beta_{j}T(\mathbf{w}_{j})$.

Second, that set is linearly independent. Suppose $\sum_{j}\beta_{j}T(\mathbf{w}_{j})=\mathbf{0}$.
Then $T\bigl(\sum_{j}\beta_{j}\mathbf{w}_{j}\bigr)=\mathbf{0}$, so
$\sum_{j}\beta_{j}\mathbf{w}_{j}\in\ker T$, and there are scalars $\alpha_{i}$ with
$\sum_{j}\beta_{j}\mathbf{w}_{j}=\sum_{i}\alpha_{i}\mathbf{u}_{i}$. Move all terms to
one side:
\[
  \sum_{j}\beta_{j}\mathbf{w}_{j}-\sum_{i}\alpha_{i}\mathbf{u}_{i}=\mathbf{0}.
\]
That relation lies inside a basis of $V$, so all its coefficients are zero; in particular every
$\beta_{j}=0$. Hence $\{T(\mathbf{w}_{j})\}$ is a basis of $\operatorname{im}T$, and
$\operatorname{rank}(T)=n-k$.
\end{proof}

Figure~\ref{fig:batang} arranges that theorem as two bars. For a space of dimension
$4$ the total length is always $4$: the identity gives $4+0$, the projection that discards one
coordinate
gives $3+1$, and the feature map $G$ later gives $2+2$. Only the position of the divider
shifts.

\begin{figure}[htbp]
\centering
\begin{tikzpicture}[x=1cm,y=1cm,font=\scriptsize,>=stealth]
  \fill[layerjade!45] (0,2.4) rectangle (4,3.1);
  \draw[layergrey] (0,2.4) rectangle (4,3.1);
  \node at (2,2.75) {rank $4$};
  \node[layerjade,anchor=east] at (-0.12,2.75) {identity};
  \node[layerjade,anchor=west] at (4.12,2.75) {nullity $0$};
  \fill[layerteal!45] (0,1.2) rectangle (3,1.9);
  \fill[layerorange!45] (3,1.2) rectangle (4,1.9);
  \draw[layergrey] (0,1.2) rectangle (4,1.9);
  \draw[layergrey] (3,1.2)--(3,1.9);
  \node at (1.5,1.55) {rank $3$};
  \node at (3.5,1.55) {$1$};
  \node[layerteal,anchor=east] at (-0.12,1.55) {projection};
  \node[layerorange,anchor=west] at (4.12,1.55) {nullity $1$};
  \fill[layerteal!45] (0,0) rectangle (2,0.7);
  \fill[layerorange!45] (2,0) rectangle (4,0.7);
  \draw[layergrey] (0,0) rectangle (4,0.7);
  \draw[layergrey] (2,0)--(2,0.7);
  \node at (1,0.35) {rank $2$};
  \node at (3,0.35) {$2$};
  \node[layerteal,anchor=east] at (-0.12,0.35) {feature map $G$};
  \node[layerorange,anchor=west] at (4.12,0.35) {nullity $2$};
  \draw[layergrey,dashed] (4,-0.25)--(4,3.4);
  \node[layergrey,anchor=south] at (4,3.4) {$\dim V=4$};
  \fill[layerteal!45] (0,-0.95) rectangle (0.35,-0.6);
  \draw[layergrey] (0,-0.95) rectangle (0.35,-0.6);
  \node[anchor=west] at (0.45,-0.78) {directions that survive};
  \fill[layerorange!45] (3.3,-0.95) rectangle (3.65,-0.6);
  \draw[layergrey] (3.3,-0.95) rectangle (3.65,-0.6);
  \node[anchor=west] at (3.75,-0.78) {directions that vanish};
\end{tikzpicture}
\caption{Rank-nullity as two bars whose sum is fixed. For $\dim V=4$: the identity map
gives $4+0$, the projection that discards one coordinate gives $3+1$, and the feature map $G$
gives
$2+2$. Only the position of the divider shifts; the total length stays $4$.}
\label{fig:batang}
\end{figure}

Two corollaries that are used most often.

\begin{akibat}
Let $\dim V=\dim W$ be finite and let $T\colon V\to W$ be linear. Then $T$ is injective if and
only if $T$ is surjective. Moreover
$\operatorname{rank}(T)\le\min(\dim V,\dim W)$.
\end{akibat}

\begin{proof}
By rank-nullity, $\dim\ker T=n-\operatorname{rank}T$ with $n=\dim V$. Injectivity holds exactly
when $\dim\ker T=0$, that is $\operatorname{rank}T=n=\dim W$, which means
$\operatorname{im}T=W$ because a subspace of full dimension fills the whole space. The
reverse implication is obtained by reversing the order of the argument. The upper bound on the rank is clear from
$\operatorname{im}T\subseteq W$ and from rank-nullity.
\end{proof}

As a numerical illustration, take $T\colon\mathbb{R}^{3}\to\mathbb{R}^{2}$ with
$T(x,y,z)=(x-y,\,y-z)$. Because $T(t,t,t)=(0,0)$, its kernel is the line $\{(t,t,t)\}$ and its nullity is
$1$.
The images $T(1,0,0)=(1,0)$ and $T(0,0,1)=(0,-1)$ span $\mathbb{R}^{2}$, so the rank is $2$;
rank-nullity: $3=1+2$.

For machine learning, that corollary is harsh. A feature map $\mathbb{R}^{n}\to\mathbb{R}^{m}$
with $n>m$ cannot possibly be injective, and its kernel has dimension at least $n-m$. Every dimensional
compression stores blind directions; the question is not whether such a direction exists, but
whether an attacker can find it.

The standard procedure for a basis of the kernel: bring the matrix to reduced row echelon form, mark the columns
without a pivot as free variables, then in turn give each free variable the value one and the other
free variables zero. The pivot variables are read off from the rows. The result is one vector per
column without a pivot, linearly independent because each has a one in a position no other
one has.

For the range, the procedure is different and is often confused with the kernel procedure above. The following
four steps are worth memorizing.

\begin{resep}
\begin{enumerate}
  \item Assemble the matrix $A$, then reduce it to reduced row echelon form by elimination.
  \item Mark the columns that contain a pivot. A column without a pivot is not a candidate basis member.
  \item Go back to the original matrix $A$, not to its echelon form. Take the original columns whose positions
        coincide with the pivot columns.
  \item Check: the number of selected vectors equals the number of nonzero rows, and those vectors are
        linearly independent.
\end{enumerate}
A warning about order: step three must not take columns from the echelon form. Row operations
change the columns, whereas the column space does not change. Taking columns from the echelon matrix
gives a wrong answer even when the count is right.
\end{resep}

The kernel also answers a more general question: once a value has been attained, all its
preimages are neatly arranged. Here $T^{-1}(\mathbf{w})$ denotes the set of preimages, not the inverse
map.

\begin{proposisi}
Let $T\colon V\to W$ be linear and $T(\mathbf{v}_{0})=\mathbf{w}$. Then
$T^{-1}(\mathbf{w})=\{\mathbf{v}_{0}+\mathbf{k}:\mathbf{k}\in\ker T\}$.
\end{proposisi}

\begin{proof}
If $\mathbf{k}\in\ker T$, then $T(\mathbf{v}_{0}+\mathbf{k})=\mathbf{w}+\mathbf{0}=\mathbf{w}$,
so every $\mathbf{v}_{0}+\mathbf{k}$ belongs to the preimage of $\mathbf{w}$. Conversely, if
$T(\mathbf{v})=\mathbf{w}$, then $T(\mathbf{v}-\mathbf{v}_{0})=\mathbf{0}$, so
$\mathbf{v}-\mathbf{v}_{0}\in\ker T$ and $\mathbf{v}=\mathbf{v}_{0}+(\mathbf{v}-\mathbf{v}_{0})$.
\end{proof}

The set of inputs that produce the same features is a shifted copy of the kernel: when the kernel
has dimension one, the family is a line, namely the blind direction of the chapter opening; when it has dimension
$k$,
an attacker has $k$ directions.

\section{Representation matrices and composition}
\tingkat{intermediate}

Now the two earlier tools are combined: coordinates turn a vector into a column, and a linear
map is determined by its values on a basis. The combination produces a matrix, where all the
computation takes place.

The representation matrix stores the image of each basis vector as one column, written
in the coordinates of the target basis.

\begin{definisi}[Representation matrix]
Let $B=(\mathbf{b}_{1},\dots,\mathbf{b}_{n})$ be a basis of $V$,
$C=(\mathbf{c}_{1},\dots,\mathbf{c}_{m})$ a basis of $W$, and $T\colon V\to W$ linear.
The \emph{representation matrix} of $T$ with respect to $B$ and $C$ is the $m\times n$ matrix
\[
[T]_{C\leftarrow
B}=\Bigl(\,[T(\mathbf{b}_{1})]_{C}\;\Big|\;\cdots\;\Big|\;[T(\mathbf{b}_{n})]_{C}\,\Bigr).
\]
\end{definisi}

Those columns are obtained by solving systems, not by guessing: for each $j$,
solve $\sum_{i}x_{i}\mathbf{c}_{i}=T(\mathbf{b}_{j})$, then write the solution as
column
$j$. When $C$ is the standard basis, that solving need not be carried out, because coordinates
with respect to the standard basis are the vector itself.

The following claim explains why the matrix is useful: it works in the expected way.

\begin{teorema}
For every $\mathbf{v}\in V$ we have
$[T(\mathbf{v})]_{C}=[T]_{C\leftarrow B}\,[\mathbf{v}]_{B}$.
\end{teorema}

\begin{proof}
Write $\mathbf{v}=\sum_{j}a_{j}\mathbf{b}_{j}$, so $[\mathbf{v}]_{B}=(a_{j})$. Linearity of
$T$ gives $T(\mathbf{v})=\sum_{j}a_{j}T(\mathbf{b}_{j})$. Taking coordinates is also
linear,
so $[T(\mathbf{v})]_{C}=\sum_{j}a_{j}[T(\mathbf{b}_{j})]_{C}$, which is exactly the matrix
product with the column $(a_{j})$.
\end{proof}

Two immediate corollaries. A linear map and its matrix determine each other, because the values on a
basis already suffice. And the rank does not change under that translation.

\begin{proposisi}
$\operatorname{rank}(T)$ equals the rank of the matrix $[T]_{C\leftarrow B}$, and $\ker T$
corresponds to the null space of that matrix. The rank of a matrix here means the dimension of its column
space, as already used in Chapter 6.
\end{proposisi}

\begin{proof}
Column $j$ of the matrix is the coordinate vector of $T(\mathbf{b}_{j})$, and
$\{T(\mathbf{b}_{j})\}$ spans $\operatorname{im}T$. Taking coordinates preserves
linear combinations, so the dimension of the range equals the dimension of the column space. The previous theorem
shows that $T(\mathbf{v})=\mathbf{0}$ exactly when $[T(\mathbf{v})]_{C}=[T]_{C\leftarrow B}[\mathbf{v}]_{B}=\mathbf{0}$,
that is, when $[\mathbf{v}]_{B}$ lies in the null space of that matrix.
\end{proof}

Composition of maps corresponds to matrix multiplication. That is no coincidence; it is
the reason matrix multiplication is defined the way it is.

\begin{teorema}
Let $T\colon V\to W$ and $S\colon W\to U$ be linear, with bases $B,C,E$ on the three spaces.
Then $[S\circ T]_{E\leftarrow B}=[S]_{E\leftarrow C}\,[T]_{C\leftarrow B}$.
\end{teorema}

\begin{proof}
For every $\mathbf{v}\in V$,
\[
  [(S\circ T)(\mathbf{v})]_{E}
  =[S]_{E\leftarrow C}[T(\mathbf{v})]_{C}
  =[S]_{E\leftarrow C}[T]_{C\leftarrow B}[\mathbf{v}]_{B}.
\]
So the matrix $[S]_{E\leftarrow C}[T]_{C\leftarrow B}$ satisfies the definition of the representation matrix
for $S\circ T$. The representation matrix is unique: if two matrices give the same result on
all $\mathbf{v}$, their results agree in particular at $\mathbf{v}=\mathbf{b}_{j}$, so
their columns agree.
\end{proof}

Composition also provides a test for invertibility through matrices.

\begin{proposisi}
Let $B$ and $C$ be bases of a space $V$ of dimension $n$, and let $T\colon V\to V$ be linear. Then
$T$ is invertible if and only if $[T]_{C\leftarrow B}$ is invertible, and in that case
$[T^{-1}]_{B\leftarrow C}=([T]_{C\leftarrow B})^{-1}$.
\end{proposisi}

\begin{proof}
If $T$ is invertible, then $T^{-1}\circ T=I$ and $T\circ T^{-1}=I$. The composition theorem gives
$[T^{-1}]_{B\leftarrow C}[T]_{C\leftarrow B}=I$ and
$[T]_{C\leftarrow B}[T^{-1}]_{B\leftarrow C}=I$, so the two matrices are mutually inverse.

Conversely, suppose $A=[T]_{C\leftarrow B}$ is invertible. If $T(\mathbf{v})=\mathbf{0}$, then
$A[\mathbf{v}]_{B}=\mathbf{0}$, so $[\mathbf{v}]_{B}=\mathbf{0}$ and
$\mathbf{v}=\mathbf{0}$.
Hence $T$ is injective. Because the source and target spaces both have dimension $n$, the
rank-nullity corollary makes $T$ surjective, so $T$ is invertible. The equation for the matrix of
$T^{-1}$
then follows from the first part.
\end{proof}

With that, a question about the inverse of an operator becomes a question about the inverse of a
matrix; the complete answer is in Chapter 8.

The special case used most often is an operator on a single space, $T\colon V\to V$,
with two different bases. Here the two matrices are linked by the transition matrix.

\begin{teorema}[Change of basis]
Let $T\colon V\to V$ be linear, $B$ and $B'$ two bases of $V$, and $P=P_{B\leftarrow B'}$.
Then
\[
  [T]_{B'}=P^{-1}[T]_{B}P .
\]
\end{teorema}

\begin{proof}
Write $T=I\circ T\circ I$ with $I$ the identity map. By the composition theorem,
\[
  [T]_{B'\leftarrow B'}=[I]_{B'\leftarrow B}\,[T]_{B\leftarrow B}\,[I]_{B\leftarrow B'}.
\]
The identity matrix with respect to two bases is the transition matrix, so
$[I]_{B'\leftarrow B}=P^{-1}$ and $[I]_{B\leftarrow B'}=P$.
\end{proof}\index{representation matrix}

Two square matrices linked in that way describe the same operator from
different viewpoints, so they are given a name of their own.

\begin{definisi}[Similar matrices]
Matrices $A$ and $B$ of size $n\times n$ are called \emph{similar} if there is an invertible
matrix $P$ with $B=P^{-1}AP$.\index{similarity}
\end{definisi}

Similarity is therefore another name for the same map seen from different bases.

\begin{proposisi}
Two $n\times n$ matrices are similar if and only if they are the representation matrices of
one and the same linear operator with respect to two bases.
\end{proposisi}

\begin{proof}
If $B=P^{-1}AP$, take the standard basis $B_{0}$ and the basis $B'$ whose columns are
the columns of $P$. Then $P=P_{B_{0}\leftarrow B'}$, and the change of basis theorem gives
$B=[T]_{B'}$ for the operator $T(\mathbf{x})=A\mathbf{x}$. Conversely, the same theorem
gives
that relation for every pair of bases.
\end{proof}

Rank is among the quantities that do not change: the rank of $[T]_{B'}$ equals the rank of $T$,
and likewise the rank of $[T]_{B}$. Determinant and trace are also invariant, but both are only
defined in Chapter 8 and discussed in Chapter 9.

A small example for the derivative operator: $D\colon P_{2}\to P_{2}$, $D(p)=p'$, with basis
$B=(1,x,x^{2})$. Because $D(1)=0$, $D(x)=1$, and $D(x^{2})=2x$,
\[
  [D]_{B}=\begin{pmatrix}0&1&0\\ 0&0&2\\ 0&0&0\end{pmatrix}.
\]
For $p=x^{2}+3x$ the coordinates are $(0,3,1)^{\top}$, and the matrix gives $(3,\,2,\,0)^{\top}$,
namely $3+2x$, exactly the derivative of $p$.

\section{Isomorphisms, dimensional invariance, and spaces of linear maps}
\tingkat{advanced}

Two different spaces can behave the same way. $P_{2}$ and $\mathbb{R}^{3}$ contain objects
of different kinds, but both have three independent directions and addition rules of the same
form.

A correspondence that preserves the whole structure is called an isomorphism: bijective,
linear, and therefore preserving bases, dimension, and kernels.

\kamus{isomorphism}{isomorfisme}{two floor plans of a house whose drawings differ but whose
number of rooms and connections between rooms are the same}

\begin{definisi}[Isomorphism]
A linear map $T\colon V\to W$ is called an \emph{isomorphism} if $T$ is bijective. Two spaces are called
\emph{isomorphic} if there is an isomorphism between them, written $V\cong W$.
\end{definisi}

One lemma before the main theorem: an isomorphism carries bases to bases.

\begin{lemma}
If $T\colon V\to W$ is an isomorphism and $B$ is a basis of $V$, then $T(B)$ is a basis of $W$.
\end{lemma}

\begin{proof}
Let $B=(\mathbf{b}_{1},\dots,\mathbf{b}_{n})$. For linear independence, suppose
$\sum_{j}c_{j}T(\mathbf{b}_{j})=\mathbf{0}$. Because $T$ is linear, the vector
$\sum_{j}c_{j}\mathbf{b}_{j}$ is mapped to $\mathbf{0}$, so that vector lies in the kernel. The kernel
is only $\{\mathbf{0}\}$ because $T$ is injective, so all $c_{j}=0$.
For spanning, take $\mathbf{w}\in W$. There is $\mathbf{v}\in V$ with $T(\mathbf{v})=\mathbf{w}$
because $T$ is surjective, and $\mathbf{v}=\sum_{j}a_{j}\mathbf{b}_{j}$ because $B$ spans.
Then
$\mathbf{w}=\sum_{j}a_{j}T(\mathbf{b}_{j})$.
\end{proof}

With that lemma, equality of dimension turns out to be the only obstacle.

\begin{teorema}
Two finite-dimensional spaces over $\mathbb{R}$ are isomorphic if and only if their dimensions are equal.
\end{teorema}

\begin{proof}
If $V\cong W$ via $T$, the previous lemma turns a basis of $V$ into a basis of $W$ with
the same number of members, so $\dim V=\dim W$.

Conversely, suppose $\dim V=\dim W=n$. Choose a basis $B$ of $V$ and $C$ of $W$. The coordinate
maps $\varphi_{B}\colon V\to\mathbb{R}^{n}$ and $\varphi_{C}\colon W\to\mathbb{R}^{n}$ are
bijective and linear. Then $\varphi_{C}^{-1}\circ\varphi_{B}$ is a bijective linear map
from
$V$ to $W$, because the composition of bijective linear maps is also bijective and linear.
\end{proof}

That conclusion is felt at once in the spaces we know.

\begin{akibat}
Every space of dimension $n$ is isomorphic to $\mathbb{R}^{n}$. In particular
$P_{n-1}\cong\mathbb{R}^{n}$ and $M_{2\times2}(\mathbb{R})\cong\mathbb{R}^{4}$.
\end{akibat}

An isomorphism preserves everything defined from the operations: linear independence, span,
kernel, range, rank, and nullity. Two isomorphic spaces cannot be distinguished
by the tools of this chapter; dimension becomes the only number that needs computing in order to recognize
a finite-dimensional space.

There is a trap in interpretation. The isomorphism above depends on the choice of
basis; there is no canonical isomorphism between $V$ and $\mathbb{R}^{n}$. Hence ``$V$ equals
$\mathbb{R}^{n}$'' is false, while ``$V$ is isomorphic to $\mathbb{R}^{n}$'' is
true.
Polynomials of degree less than three and triples of real numbers are objects of different kinds,
but the two are isomorphic because every polynomial is determined by its three coefficients.

Now the collection of all linear maps between two spaces; that set is itself a vector space.

\begin{teorema}
The set $\mathcal{L}(V,W)$ of all linear maps $V\to W$, with the operations
$(S+T)(\mathbf{v})=S(\mathbf{v})+T(\mathbf{v})$ and $(cT)(\mathbf{v})=cT(\mathbf{v})$,
is
a vector space. If $\dim V=n$ and $\dim W=m$, then $\dim\mathcal{L}(V,W)=mn$. The mapping
$T\mapsto[T]_{C\leftarrow B}$ is an isomorphism $\mathcal{L}(V,W)\to M_{m\times n}(\mathbb{R})$.
\end{teorema}

\begin{proof}
$\mathcal{L}(V,W)$ sits inside the space of all functions $V\to W$, which is a vector space
by the example in Chapter 6. The zero map is linear. The sum of two linear maps is linear, and a scalar multiple of a
linear
map is also linear; so $\mathcal{L}(V,W)$ is a subspace, and therefore a vector space.

For the isomorphism claim, let $\Phi(T)=[T]_{C\leftarrow B}$. Linearity of $\Phi$: column
$j$ of $\Phi(S+cT)$ is $[S(\mathbf{b}_{j})+cT(\mathbf{b}_{j})]_{C}$, which equals
$[S(\mathbf{b}_{j})]_{C}+c[T(\mathbf{b}_{j})]_{C}$ because coordinates are linear. Hence
$\Phi(S+cT)=\Phi(S)+c\Phi(T)$.

Injectivity: if $\Phi(T)$ is the zero matrix, then all $[T(\mathbf{b}_{j})]_{C}=\mathbf{0}$, so
$T(\mathbf{b}_{j})=\mathbf{0}$ for all $j$, and $T=\mathbf{0}$ because a linear map is
determined by its values on a basis.

Surjectivity: given a matrix $A=(a_{ij})$ of size $m\times n$. Define
$T(\mathbf{b}_{j})=\sum_{i}a_{ij}\mathbf{c}_{i}$ and extend linearly; by the
proposition
on values on a basis, that extension exists and is unique. Column $j$ of $\Phi(T)$ is
$(a_{ij})_{i}$, so $\Phi(T)=A$.

Because $\Phi$ is an isomorphism and $\dim M_{m\times n}(\mathbb{R})=mn$, we get
$\dim\mathcal{L}(V,W)=mn$.
\end{proof}

That count gives an immediate corollary.

\begin{akibat}
$\dim\mathcal{L}(V,V)=n^{2}$ if $\dim V=n$. For example,
$\dim\mathcal{L}(P_{2},P_{2})=9$ and $\dim\mathcal{L}(P_{2},\mathbb{R}^{2})=6$.
\end{akibat}

That last result explains why a matrix is more than a computational device. The space of operators
and the space of matrices have the same dimension and are linked by an isomorphism that respects composition,
so every statement about operators has a counterpart about matrices. What is not
inherited
is simplicity of form: the same operator can be complicated in one basis and simple
in
another; the search for the best basis is the program of Chapter 9.

\begin{tantangan}
Find one linear map $T\colon\mathbb{R}^{3}\to\mathbb{R}^{3}$ that maps the unit square
in the $xy$ plane to another unit square of area still $1$, yet
$\operatorname{rank}(T)=2$. Explain why an area preserved in one plane does not
force
full rank.

Hint 1: examine one plane, not the whole space. After that, ask where the third axis
goes.

Hint 2: place a $90^{\circ}$ rotation on the first two coordinates and send the third coordinate to zero,
so that $T(\mathbf{e}_{3})=\mathbf{0}$. Compute the rank from the images of $\mathbf{e}_{1}$ and
$\mathbf{e}_{2}$.

If you get stuck, return to the section Linear transformations, kernel, and range, especially the definition of
rank and the rank-nullity corollary.
\end{tantangan}

\begin{untukanak}
Two classes at school can differ in size, yet have the same number of chairs. If the
number is
the same, every pupil in the first class can be paired with one chair in the second, with no
empty
chair and no pupil getting two chairs. As far as counting chairs goes, the two classes are
equivalent.

What is not equivalent is the pupils' names and faces. The chair pairing does not say
who sits where; it only pairs. Different people can make different pairings,
and all of them are legitimate.

In this chapter we will learn when two spaces that look different may be regarded as the same for
computational purposes, and when that assumption is misleading. The equivalence always depends
on
our choice, not on something already present in the spaces.

If two classes have the same number of chairs, may the headmaster simply swap the class names
just
like that? The answer is no, and the reason is in this chapter.
\end{untukanak}

\section{A worked example}
\tingkat{intermediate}

Consider the matrix
\[
  A=\begin{pmatrix}1&2&1\\ 2&5&3\\ 3&7&4\end{pmatrix},
  \qquad \text{third row}=\text{first row}+\text{second row}.
\]
We determine its rank, its nullity, and both bases, then examine the same operator in another
basis.\index{rank}\index{nullity}

Row elimination:
\[
  \begin{pmatrix}1&2&1\\ 2&5&3\\ 3&7&4\end{pmatrix}
  \xrightarrow[R_{3}\leftarrow R_{3}-3R_{1}]{R_{2}\leftarrow R_{2}-2R_{1}}
  \begin{pmatrix}1&2&1\\ 0&1&1\\ 0&1&1\end{pmatrix}
  \xrightarrow{R_{3}\leftarrow R_{3}-R_{2}}
  \begin{pmatrix}1&2&1\\ 0&1&1\\ 0&0&0\end{pmatrix}.
\]
Then $R_{1}\leftarrow R_{1}-2R_{2}$ gives the reduced echelon form
\[
  \begin{pmatrix}1&0&-1\\ 0&1&1\\ 0&0&0\end{pmatrix}.
\]
There are two nonzero rows, so $\operatorname{rank}A=2$ and $\det A=0$; the pivot columns are the
first and the second.

Kernel: from $x_{1}-x_{3}=0$ and $x_{2}+x_{3}=0$ we obtain $x_{1}=x_{3}$ and
$x_{2}=-x_{3}$.
Hence
\[
  \ker A=\operatorname{span}\{(1,-1,1)^{\top}\},\qquad \operatorname{nullity}(A)=1 .
\]
Rank-nullity is satisfied: $3=2+1$.

Range: the columns of $A$ are $\mathbf{c}_{1}=(1,2,3)^{\top}$,
$\mathbf{c}_{2}=(2,5,7)^{\top}$, and $\mathbf{c}_{3}=(1,3,4)^{\top}$. The first and the
second
pivot columns give the basis $\{\mathbf{c}_{1},\mathbf{c}_{2}\}$. The third column is not independent:
$\mathbf{c}_{2}-\mathbf{c}_{1}=\mathbf{c}_{3}$, and that relation is exactly the kernel vector
$(1,-1,1)$: $\mathbf{c}_{1}-\mathbf{c}_{2}+\mathbf{c}_{3}=\mathbf{0}$.

Now consider $T(\mathbf{x})=A\mathbf{x}$ and the basis
$B'=(\mathbf{b}'_{1},\mathbf{b}'_{2},\mathbf{b}'_{3})$ with
$\mathbf{b}'_{1}=(1,1,0)$, $\mathbf{b}'_{2}=(1,0,1)$, $\mathbf{b}'_{3}=(0,1,1)$. Its
transition matrix from the standard basis $B$ is
\[
  P=P_{B\leftarrow B'}=\begin{pmatrix}1&1&0\\ 1&0&1\\ 0&1&1\end{pmatrix},
  \qquad
  P^{-1}=\frac{1}{2}\begin{pmatrix}1&1&-1\\ 1&-1&1\\ -1&1&1\end{pmatrix}.
\]
The images of the new basis vectors are computed first:
\[
  A\mathbf{b}'_{1}=\begin{pmatrix}3\\7\\10\end{pmatrix},\quad
  A\mathbf{b}'_{2}=\begin{pmatrix}2\\5\\7\end{pmatrix},\quad
  A\mathbf{b}'_{3}=\begin{pmatrix}3\\8\\11\end{pmatrix}.
\]
Hence
\[
  AP=\begin{pmatrix}3&2&3\\ 7&5&8\\ 10&7&11\end{pmatrix},
  \qquad
  [T]_{B'}=P^{-1}AP=\begin{pmatrix}0&0&0\\ 3&2&3\\ 7&5&8\end{pmatrix}.
\]
Check one entry: first row, first column $\tfrac{1}{2}(3+7-10)=0$; second row,
third column $\tfrac{1}{2}(3-8+11)=3$. The rest follows the same pattern.

A first row that is entirely zero is no accident: all three images lie inside
$\operatorname{span}(\mathbf{b}'_{2},\mathbf{b}'_{3})$, for instance
$A\mathbf{b}'_{3}=(3,8,11)=3\mathbf{b}'_{2}+8\mathbf{b}'_{3}$. Hence the first coordinate of
every image is zero, and the range of $T$ is contained in the coordinate plane of the first
coordinate.

The pattern is general. A zero row signals a range contained in a coordinate subspace
of lower dimension; a zero column signals a basis vector that goes into the kernel. Both symptoms
are visible only in the chosen basis, whereas rank and nullity do not change.

A complete check:
\[
  \operatorname{rank}[T]_{B'}=2,\qquad \operatorname{trace}[T]_{B'}=0+2+8=10=\operatorname{trace}A,
  \qquad \det[T]_{B'}=0=\det A .
\]
The equality of traces is not an arithmetic accident. The $j$-th diagonal entry is the coefficient of
$\mathbf{b}'_{j}$ in the image $T(\mathbf{b}'_{j})$, that is, how much $T$
preserves
that direction itself. The sum does not depend on the directions chosen, only on the operator,
so the trace serves as a cheap fingerprint: different traces mean the matrices are not similar, equal
traces
guarantee nothing.

For $\mathbf{v}=(3,8,11)$ we have $[\mathbf{v}]_{B'}=(0,3,8)^{\top}$ and
\[
  [T]_{B'}[\mathbf{v}]_{B'}=\begin{pmatrix}0\\30\\79\end{pmatrix},
\]
whereas $T(\mathbf{v})=A\mathbf{v}=(30,79,109)^{\top}$ and its coordinates are
$(0,30,79)^{\top}$. The two agree.

The feature map $G\colon\mathbb{R}^{4}\to\mathbb{R}^{3}$
is given by the matrix
\[
  G=\begin{pmatrix}1&2&0&1\\ 0&1&1&1\\ 1&3&1&2\end{pmatrix},
\]
whose third row is the sum of the first two rows. Elimination gives the reduced echelon
form
\[
  \begin{pmatrix}1&0&-2&-1\\ 0&1&1&1\\ 0&0&0&0\end{pmatrix},
\]
so $\operatorname{rank}G=2$ and $\operatorname{nullity}(G)=4-2=2$. From
$x_{1}=2x_{2}+x_{4}$ and $x_{3}=-x_{2}-x_{4}$:
\[
  \mathbf{k}_{1}=(-2,1,-1,0)^{\top},\qquad \mathbf{k}_{2}=(-1,0,-1,1)^{\top}.
\]
Both satisfy $G\mathbf{k}_{i}=\mathbf{0}$ and are linearly independent, so the changes that the
three features do not see form a two-dimensional plane, not a single line: an attacker can
shift the input in any direction of that plane without moving the score.

\kertas{Repeat the computation with the basis $B''=(\mathbf{b}'_{2},\mathbf{b}'_{3},\mathbf{b}'_{1})$,
that is, with $B'$ permuted. Predict how the zero row and the other entries move, then compute
$[T]_{B''}$ via $Q^{-1}[T]_{B'}Q$ with $Q$ the corresponding permutation matrix. Check that
rank
and trace do not change, and explain why $\det A=0$ can be concluded directly from ``the third
row is the sum of the first two rows''.}

All the numbers above were checked with exact arithmetic (SymPy); the numpy and
scipy checks are in the next section, including one number that is not exactly zero because of machine
rounding.

\section{Computational lab}
\tingkat{advanced}

\begin{lab}
\begin{verbatim}
import numpy as np
from scipy.linalg import null_space

A = np.array([[1., 2., 1.],
              [2., 5., 3.],
              [3., 7., 4.]])
print("rank A  :", np.linalg.matrix_rank(A))
print("det A   :", np.linalg.det(A))
K = null_space(A)
print("dim ker :", K.shape[1])
print("basis ker:\n", np.round(K, 4))
print("A @ K   :\n", np.round(A @ K, 12))

P = np.array([[1., 1., 0.],
              [1., 0., 1.],
              [0., 1., 1.]])
B = np.linalg.inv(P) @ A @ P
print("[T]_B'  :\n", np.round(B, 10))
print("difference :", np.linalg.norm(B - np.linalg.inv(P) @ A @ P))
print("trace   :", round(np.trace(B), 10), round(np.trace(A), 10))
print("rank B  :", np.linalg.matrix_rank(B))
print("cond P  :", np.linalg.cond(P))

v = np.array([3., 8., 11.])
print("[v]_B'  :", np.round(np.linalg.solve(P, v), 4))
print("[Av]_B' :", np.round(np.linalg.solve(P, A @ v), 4))
print("B[v]_B' :", np.round(B @ np.linalg.solve(P, v), 4))

x0 = np.array([1., 0., 0.])
d = 0.5 * K[:, 0]
print("A@d     :", np.round(A @ d, 12))
print("shift   :", np.linalg.norm(A @ (x0 + d) - A @ x0))
\end{verbatim}
Actual output:
\begin{verbatim}
rank A  : 2
det A   : 6.66133814775094e-16
dim ker : 1
basis ker:
 [[ 0.5774]
 [-0.5774]
 [ 0.5774]]
A @ K   :
 [[-0.]
 [ 0.]
 [ 0.]]
[T]_B'  :
 [[0. 0. 0.]
 [3. 2. 3.]
 [7. 5. 8.]]
difference : 0.0
trace   : 10.0 10.0
rank B  : 2
cond P  : 2.0000000000000004
[v]_B'  : [0. 3. 8.]
[Av]_B' : [ 0. 30. 79.]
B[v]_B' : [ 0. 30. 79.]
A@d     : [-0.  0.  0.]
shift   : 6.280369834735101e-16
\end{verbatim}
\end{lab}

Four things are worth reading from that output.

First, $\det A$ is not zero but $6.66\times10^{-16}$. This matrix is singular
exactly, but binary fraction arithmetic knows no exact zero.
\texttt{matrix\_rank} compares the largest and smallest singular values, not a simple threshold,
so the rank is still reported as $2$; Chapter 22 discusses why such a threshold is necessary.

The choice of singular values rather than the determinant: the determinant measures volume, and a volume can be
very
small without the rank changing. What determines the rank is the ratio of the smallest singular value
to the
largest.

Second, \texttt{null\_space} returns a unit vector with entries $0.5774$,
$-0.5774$, and $0.5774$, that is $(1,-1,1)^{\top}$ divided by $\sqrt{3}$, and
$1/\sqrt{3}=0.57735$ agrees. The direction is the same as in the hand computation; its length is normalized
because
scipy chooses an orthonormal basis. A kernel basis is not unique, its dimension is.

Third, the $[T]_{B'}$ reported by numpy matches the hand computation exactly, and the traces
of both matrices agree, $10$. The difference of $P^{-1}AP$ computed twice is $0.0$ because
$P$
contains small integers. The condition number of $P$ is $2.0000000000000004$, quite
small;
a bad basis inflates the error, and that is discussed in Chapter 10.

Fourth, the kernel direction is genuinely blind to $T$. Pushing $\mathbf{x}_{0}=(1,0,0)$ by
half the kernel vector changes $\mathbf{x}_{0}$ visibly, yet $A\mathbf{x}$ shifts only
$6.28\times10^{-16}$, that is, zero at machine precision.

A final warning concerns the choice of basis. A nearly dependent basis makes coordinates
explode, and the reported numbers lose their meaning.

\begin{lab}
\begin{verbatim}
Pb = np.array([[1., 1., 0.],
               [1., 1.001, 0.],
               [0., 0., 1.]])
print("cond Pb:", np.linalg.cond(Pb))
print("coord   :", np.linalg.inv(Pb) @ v)
print("error   :", np.linalg.norm(np.linalg.inv(Pb) @ (Pb @ v) - v))
\end{verbatim}
Actual output:
\begin{verbatim}
cond Pb: 4002.00075012484
coord   : [-4997.  5000.    11.]
error   : 4.823324033076557e-13
\end{verbatim}
\end{lab}

The first two basis vectors differ only in the third digit, and the condition number of $P_{b}$
reaches $4002$. The coordinates $(3,8,11)$ explode into $(-4997,\,5000,\,11)$. The reconstruction
error
is still $4.8\times10^{-13}$, but the relative error in the coordinates grows by a factor of
the
condition number; with 16 digits of precision, about 12 remain. In a worse basis that remainder
is
used up and the reported rank cannot be trusted.

\begin{untukdosen}
A teaching order that works in practice: coordinates first, linear maps second. Students who
know
coordinates before the abstract definition are readier to accept the representation matrix, because
the matrix appears as a list of coordinates, not as a new object. After rank-nullity, move to
change of basis.

Two typical misconceptions. First, equating a vector with its coordinate list, then getting confused
when two bases give two lists for the same vector. Second, the direction of the transition matrix
backwards. A quick check: the columns must contain the addresses of the basis vectors, and the subscript is read
right to left.

Blackboard activity: give one operator on $\mathbb{R}^{2}$ and two bases. Ask half the
class to compute $[T]_{B'}$ via the column definition, half via $P^{-1}AP$, then compare
on
the board; the noise that arises is usually exactly about the direction of $P$.

A sign of readiness for the next step: students can explain why a zero row appears in the new
matrix, without recomputing. If they connect it with a range
contained in a coordinate subspace, they are ready for Chapter 9.
\end{untukdosen}

\section{Summary}
\tingkat{basic}

\begin{ringkasan}
\begin{itemize}
  \item The coordinates $[\mathbf{v}]_{B}$ are the unique list of coefficients with respect to an ordered basis;
        the map $\mathbf{v}\mapsto[\mathbf{v}]_{B}$ is linear and bijective, but depends on the
        basis.
  \item The transition matrix contains the coordinates of the old basis in the new basis:
        $[\mathbf{v}]_{C}=P_{C\leftarrow B}[\mathbf{v}]_{B}$, with
        $P_{B\leftarrow C}=P_{C\leftarrow B}^{-1}$.
  \item A linear map is determined by its values on a basis; $T$ is injective exactly when
        $\ker T=\{\mathbf{0}\}$.
  \item Rank-nullity: $\dim V=\dim\ker T+\dim\operatorname{im}T$; in equal dimensions,
        injective is equivalent to surjective.
  \item The representation matrix satisfies
$[T(\mathbf{v})]_{C}=[T]_{C\leftarrow B}[\mathbf{v}]_{B}$, composition becomes matrix
        multiplication, and $[T]_{B'}=P^{-1}[T]_{B}P$.
  \item $\dim\mathcal{L}(V,W)=mn$ if $\dim V=n$, $\dim W=m$; two finite-dimensional spaces
        are isomorphic exactly when their dimensions are equal.
\end{itemize}
\end{ringkasan}

\section{Exercises}
\tingkat{intermediate}

\begin{latihan}
\begin{enumerate}
  \item With $B'=(\mathbf{b}'_{1},\mathbf{b}'_{2},\mathbf{b}'_{3})$ as in the worked
example, determine $[\mathbf{w}]_{B'}$ for $\mathbf{w}=(2,-1,3)$, then check the result
        by multiplying back.

\item Prove that every isomorphism $V\to\mathbb{R}^{n}$ has the form $\varphi_{B}$ for
        some basis $B$, and explain why an isomorphism to $\mathbb{R}^{n}$ is not unique.
        (Hint: take the preimages of the standard basis, then show that this set is a basis.)

  \item For $A=\begin{pmatrix}0&1&0\\ 0&0&2\\ 0&0&0\end{pmatrix}$, determine the rank, the nullity,
        a basis of the kernel, and a basis of the range. Explain the connection with the derivative operator on
        $P_{2}$.

\item Prove $\operatorname{rank}(S\circ T)\le\min(\operatorname{rank}S,\operatorname{rank}T)$
        for linear maps that are defined. Use
        $\operatorname{im}(S\circ T)\subseteq\operatorname{im}S$ and compare the kernels.
        (Hint: for the second bound, compare $\ker T$ with $\ker(S\circ T)$ via
        rank-nullity.)

\item Assess this claim: \emph{two square matrices of the same size that have the same rank
are necessarily similar.} Decide whether it is true, false, or not supported by enough evidence, with reasons and a counterexample if needed.

\item Let $T\colon\mathbb{R}^{3}\to\mathbb{R}^{2}$ with
        $T(x,y,z)=(x+2y,\;y+z)$. Compute $[T]_{C\leftarrow B}$ for
$B=(\,(1,0,0),(1,1,0),(1,1,1)\,)$ and $C=(\,(1,1),(1,0)\,)$, then check the formula at
        $\mathbf{v}=(1,2,3)$.
        (Hint: $[T]_{C\leftarrow B}=C^{-1}AB$; compute $AB$ first, then solve
        $C\mathbf{x}=\mathbf{y}$ for each of its columns.)

  \item Prove: if $U$ is a subspace of a finite-dimensional $V$ and $\dim U=\dim V$, then $U=V$.

  \item Prove that if $A$ and $B$ are similar then $\operatorname{nullity}(A)
        =\operatorname{nullity}(B)$. Explain why that statement is equivalent to
        ``the kernel does not depend on the basis''.

  \item The feature map $F\colon\mathbb{R}^{4}\to\mathbb{R}^{3}$ is given by the matrix
        \[
          F=\begin{pmatrix}1&0&2&1\\ 0&1&1&3\\ 1&1&3&4\end{pmatrix}.
        \]
        Determine the rank and nullity of $F$, give one nonzero kernel vector, and explain why
        that vector is a direction of change that does not shift a single feature.
\end{enumerate}
\end{latihan}

\begin{untukpeneliti}
What is proved in this chapter are statements about finite dimension over the field $\mathbb{R}$:
rank-nullity, uniqueness of the representation matrix, the equivalence of isomorphism with equality of
dimension,
and $\dim\mathcal{L}(V,W)=mn$. All of them are theorems, not conventions. The claim ``a matrix is its transformation'' is merely popular; without a choice of basis that claim is meaningless, and the zero-row
example
above shows that the form of a matrix depends on the basis.

Primary sources: Axler, \emph{Linear Algebra Done Right} 4th edition, sections 3A to 3D,
contain
spaces of linear maps, kernel and range, matrices, isomorphisms, and change of basis;
rank-nullity
there is called the fundamental theorem for linear maps. Hoffman and Kunze, \emph{Linear
Algebra}
2nd edition, Chapter 3, is denser but equivalent.

An honest open question. The classification of operators up to similarity is not given
here; Chapter 9 stops at diagonalization. Numerically, the rank decision depends on a
threshold and
condition numbers; this chapter uses library output without discussing its stability, a debt paid
in Chapter 22. Limit of validity: in infinite-dimensional spaces the whole argument collapses,
rank-nullity
fails, and isomorphism is no longer equivalent to equality of dimension.
\end{untukpeneliti}

\peta{This chapter follows Anton, \emph{Elementary Linear Algebra} 11th edition, sections 4.4, 4.6,
4.8, and 8.1 to 8.5, with the abstract counterpart in Axler, \emph{Linear Algebra Done Right}
4th
edition, Chapter 3, and Hoffman and Kunze, \emph{Linear Algebra} 2nd edition, Chapter 3.}
